% =========================================================
% Intermediate Value Theorem — Notes
% Chapter: continuity
% Section: ivt
% Volume III · Analysis
% =========================================================

\section{Intermediate Value Theorem}
\label{sec:ivt}

% ---------------------------------------------------------
% Toolkit
% ---------------------------------------------------------

\begin{tcolorbox}[title={Toolkit — Boundedness, Extrema, and Intermediate Values}]
\begin{itemize}
  \item Implication chain on $[a,b]$: continuity $\Rightarrow$ bounded
    $\Rightarrow$ extrema attained $\Rightarrow$ all intermediate values
    attained $\Rightarrow$ image is a closed bounded interval.
  \item The three hypotheses are essential: closed, bounded, continuous.
    Drop any one and the conclusion fails.
  \item IVT on a general interval: continuous image of any interval is
    an interval.
  \item Darboux property: no value is skipped. The sign-change form
    (Location of Roots) is the operative tool for root-finding arguments.
\end{itemize}
\end{tcolorbox}

% ---------------------------------------------------------
% Definition — Bounded on a Set
% ---------------------------------------------------------

\begin{definition}[Bounded on a Set]
\label{def:bounded-on-a-set}
Let $A \subseteq \mathbb{R}$ and let $f : A \to \mathbb{R}$. The
function $f$ is \textbf{bounded on $A$} if there exists $M > 0$
such that $|f(x)| \leq M$ for all $x \in A$.
\end{definition}

\begin{remark*}[Standard quantified statement]
\[
  f \text{ bounded on } A
  \;\iff\;
  \exists M > 0,\;
  \forall x \in A:\;
  |f(x)| \leq M.
\]
\end{remark*}

\begin{remark*}[Negated quantified statement]
\[
  f \text{ unbounded on } A
  \;\iff\;
  \forall M > 0,\;
  \exists x \in A:\;
  |f(x)| > M.
\]
\end{remark*}

\begin{remark*}[Interpretation]
Boundedness means the image $f(A)$ fits inside the interval $[-M, M]$
for some finite $M$. The exact value of $M$ is irrelevant; what
matters is that a single global bound holds across the entire domain.
Boundedness is the minimal global control on $f$ and is the first
consequence of continuity on a compact interval.
\end{remark*}

% ---------------------------------------------------------
% Theorem — Boundedness Theorem
% ---------------------------------------------------------

\begin{theorem}[Boundedness Theorem]
\label{thm:boundedness-theorem}
Let $I = [a,b]$ be a closed bounded interval and let
$f : I \to \mathbb{R}$ be continuous on $I$. Then $f$ is bounded
on $I$.

\smallskip
\noindent
\hyperref[prf:boundedness-theorem]{\textit{Go to proof.}}
\end{theorem}

\begin{remark*}[Standard quantified statement]
\[
  \forall a, b \in \mathbb{R},\; a \leq b,\;
  \forall f : [a,b] \to \mathbb{R}:\;
  f \text{ cts on } [a,b]
  \;\Rightarrow\;
  \exists M > 0,\;
  \forall x \in [a,b]:\;
  |f(x)| \leq M.
\]
\end{remark*}

\begin{remark*}[Negated quantified statement]
\[
  \bigl(\forall M > 0,\; \exists x \in [a,b]:\; |f(x)| > M\bigr)
  \;\Rightarrow\;
  f \text{ is not continuous on } [a,b].
\]
\end{remark*}

\begin{remark*}[Interpretation]
Boundedness is the first compactness consequence for continuous
functions on closed bounded intervals. The proof proceeds by
contradiction via the Bolzano--Weierstrass theorem: if $f$ were
unbounded, a sequence $(x_n) \subseteq [a,b]$ with $|f(x_n)| > n$
would exist; compactness produces a convergent subsequence; continuity
forces $f$ to be bounded near the limit — contradiction. The three
hypotheses are each necessary: $f(x) = x$ on $(0,1)$ (not closed),
$f(x) = x$ on $[0,\infty)$ (not bounded), and $f(x) = 1/x$ on
$[0,1]$ with $f(0)$ undefined (not continuous) all fail to be bounded.
\end{remark*}

% ---------------------------------------------------------
% Definition — Absolute Maximum and Minimum
% ---------------------------------------------------------

\begin{tcolorbox}[title={Definition --- Absolute Extrema}]
\begin{definition}[Absolute Maximum and Absolute Minimum]
\label{def:absolute-extrema}
Let $A \subseteq \mathbb{R}$ and let $f : A \to \mathbb{R}$.

The function $f$ has an \textbf{absolute maximum on $A$} if there
exists $x^* \in A$ such that $f(x^*) \geq f(x)$ for all $x \in A$.

The function $f$ has an \textbf{absolute minimum on $A$} if there
exists $x_* \in A$ such that $f(x_*) \leq f(x)$ for all $x \in A$.

The points $x^*$ and $x_*$ are called \textbf{absolute maximum point}
and \textbf{absolute minimum point} of $f$ on $A$ respectively.
\end{definition}
\end{tcolorbox}

\begin{remark*}[Standard quantified statement]
\[
\begin{aligned}
  f \text{ has abs max on } A
  &\;\iff\;
  \exists x^* \in A,\;
  \forall x \in A:\;
  f(x^*) \geq f(x). \\
  f \text{ has abs min on } A
  &\;\iff\;
  \exists x_* \in A,\;
  \forall x \in A:\;
  f(x_*) \leq f(x).
\end{aligned}
\]
\end{remark*}

\begin{remark*}[Interpretation]
An absolute extremum is a global attainment: a point in the domain
where the supremum or infimum of $f$ is actually achieved. Being
bounded above implies the supremum exists, but does not guarantee it
is attained. The defining feature of an absolute maximum is that the
supremal value occurs at some $x^* \in A$, not merely as a limit.
\end{remark*}

% ---------------------------------------------------------
% Theorem — Extreme Value Theorem
% ---------------------------------------------------------

\begin{theorem}[Extreme Value Theorem]
\label{thm:extreme-value-theorem}
Let $I = [a,b]$ be a closed bounded interval and let
$f : I \to \mathbb{R}$ be continuous on $I$. Then $f$ has an
absolute maximum and an absolute minimum on $I$.

\smallskip
\noindent
\hyperref[prf:evt-maximum-minimum]{\textit{Go to proof.}}
\end{theorem}

\begin{remark*}[Standard quantified statement]
\[
  f \text{ cts on } [a,b]
  \;\Rightarrow\;
  \Bigl(\exists x^* \in [a,b],\; \forall x \in [a,b]:\; f(x^*) \geq f(x)\Bigr)
  \;\wedge\;
  \Bigl(\exists x_* \in [a,b],\; \forall x \in [a,b]:\; f(x_*) \leq f(x)\Bigr).
\]
\end{remark*}

\begin{remark*}[Interpretation]
The Extreme Value Theorem upgrades boundedness to attainment.
Boundedness guarantees the supremum and infimum exist; the EVT
guarantees they are achieved at actual points of $[a,b]$. The
proof extracts a maximizing sequence, applies Bolzano--Weierstrass
to get a convergent subsequence, and uses continuity to show the
limit point achieves the supremum. Compactness ($[a,b]$ closed and
bounded) is essential at both steps.
\end{remark*}

% ---------------------------------------------------------
% Theorem — Image of a Closed Bounded Interval Is Closed and Bounded
% ---------------------------------------------------------

\begin{theorem}[Image of a Closed Bounded Interval Is Closed and Bounded]
\label{thm:image-of-closed-bounded-interval-is-closed-bounded}
Let $I = [a,b]$ be a closed bounded interval and let
$f : I \to \mathbb{R}$ be continuous on $I$. Then $f(I)$ is closed
and bounded.

\smallskip
\noindent
\hyperref[prf:image-of-function-on-closed-bounded-interval-is-closed-bounded]{\textit{Go to proof.}}
\end{theorem}

\begin{remark*}[Standard quantified statement]
\[
  f \text{ cts on } [a,b]
  \;\Rightarrow\;
  f([a,b]) \text{ is closed and bounded.}
\]
\end{remark*}

\begin{remark*}[Interpretation]
Boundedness of $f(I)$ follows from the Boundedness Theorem.
Closedness follows because $[a,b]$ is compact and the continuous
image of a compact set is compact, and compact subsets of
$\mathbb{R}$ are closed and bounded. This theorem is the
prerequisite for the result that $f(I)$ is a closed bounded interval.
\end{remark*}

% ---------------------------------------------------------
% Theorem — Location of Roots
% ---------------------------------------------------------

\begin{theorem}[Location of Roots]
\label{thm:location-of-roots}
Let $I = [a,b]$ and let $f : I \to \mathbb{R}$ be continuous on $I$.
If $f(a) < 0 < f(b)$ or $f(a) > 0 > f(b)$, then there exists
$c \in (a,b)$ such that $f(c) = 0$.

\smallskip
\noindent
\hyperref[prf:locations-of-roots]{\textit{Go to proof.}}
\end{theorem}

\begin{remark*}[Standard quantified statement]
\[
  f \text{ cts on } [a,b]
  \;\wedge\;
  \bigl(f(a) < 0 < f(b) \;\vee\; f(a) > 0 > f(b)\bigr)
  \;\Rightarrow\;
  \exists c \in (a,b):\; f(c) = 0.
\]
\end{remark*}

\begin{remark*}[Interpretation]
A continuous function cannot pass from negative to positive values
without crossing zero. The proof uses the bisection method: maintain
a nested sequence of intervals on which the function changes sign;
the intersection point is the root. This is the sign-change form of
the intermediate value principle and the operative tool in all
root-finding arguments.
\end{remark*}

% ---------------------------------------------------------
% Theorem — Bolzano's Intermediate Value Theorem
% ---------------------------------------------------------

\begin{theorem}[Bolzano's Intermediate Value Theorem]
\label{thm:bolzano-intermediate-value}
Let $I$ be an interval and let $f : I \to \mathbb{R}$ be continuous
on $I$. If $a, b \in I$ and $k \in \mathbb{R}$ satisfies
$f(a) < k < f(b)$, then there exists $c \in I$ strictly between
$a$ and $b$ such that $f(c) = k$.

\smallskip
\noindent
\hyperref[prf:bolzano-ivt]{\textit{Go to proof.}}
\end{theorem}

\begin{remark*}[Standard quantified statement]
\[
  f \text{ cts on } I,\;
  a, b \in I,\;
  f(a) < k < f(b)
  \;\Rightarrow\;
  \exists c \in I:\;
  \min\{a,b\} < c < \max\{a,b\}
  \;\wedge\;
  f(c) = k.
\]
\end{remark*}

\begin{remark*}[Negated quantified statement]
\[
  \bigl(\forall c \in I:\; f(c) \neq k\bigr)
  \;\Rightarrow\;
  \neg\bigl(f(a) < k < f(b)\bigr)
  \;\vee\;
  f \text{ not cts on } I.
\]
\end{remark*}

\begin{remark*}[Interpretation]
The IVT is the central structural result connecting continuity and
the order structure of $\mathbb{R}$: continuous functions on
intervals do not skip intermediate values. The proof reduces to the
Location of Roots theorem by considering $g = f - k$. The
generalization to arbitrary intervals (not just $[a,b]$) is
immediate since every interval in $\mathbb{R}$ is connected. This
is the topological content: continuous images of connected sets are
connected, and the connected subsets of $\mathbb{R}$ are exactly
the intervals.
\end{remark*}

% ---------------------------------------------------------
% Corollary — Every Value Between inf and sup Is Attained
% ---------------------------------------------------------

\begin{corollary}[Every Value Between $\inf f(I)$ and $\sup f(I)$ Is Attained]
\label{cor:ivt-between-inf-and-sup}
Let $I = [a,b]$ be a closed bounded interval and let
$f : I \to \mathbb{R}$ be continuous on $I$. If $k \in \mathbb{R}$
satisfies $\inf f(I) \leq k \leq \sup f(I)$, then there exists
$c \in I$ such that $f(c) = k$.

\smallskip
\noindent
\hyperref[prf:every-value-attained]{\textit{Go to proof.}}
\end{corollary}

\begin{remark*}[Standard quantified statement]
\[
  f \text{ cts on } [a,b],\;
  \inf f([a,b]) \leq k \leq \sup f([a,b])
  \;\Rightarrow\;
  \exists c \in [a,b]:\; f(c) = k.
\]
\end{remark*}

\begin{remark*}[Interpretation]
The EVT identifies points where the infimum and supremum are
attained. Bolzano's IVT then fills in every intermediate value.
The corollary combines these: continuity on $[a,b]$ forces the image
to contain every value between its extreme values, without gaps.
\end{remark*}

% ---------------------------------------------------------
% Theorem — Image of a Closed Bounded Interval Is an Interval
% ---------------------------------------------------------

\begin{theorem}[Image of a Closed Bounded Interval Is a Closed Bounded Interval]
\label{thm:image-of-closed-bounded-interval}
Let $I = [a,b]$ be a closed bounded interval and let
$f : I \to \mathbb{R}$ be continuous on $I$. Then $f(I)$ is a
closed bounded interval.

\smallskip
\noindent
\hyperref[prf:image-of-closed-bounded-interval-is-interval]{\textit{Go to proof.}}
\end{theorem}

\begin{remark*}[Standard quantified statement]
\[
  f \text{ cts on } [a,b]
  \;\Rightarrow\;
  f([a,b]) = [\,\min f,\, \max f\,]
  \quad\text{for some }\min f, \max f \in \mathbb{R}.
\]
\end{remark*}

\begin{remark*}[Interpretation]
This theorem synthesizes the three preceding results. The image is
bounded (Boundedness Theorem), its extreme values are attained (EVT),
and every intermediate value is attained (IVT). Together they force
$f(I)$ to be precisely the closed interval $[\min f, \max f]$ — no
gaps, both endpoints included.
\end{remark*}

% ---------------------------------------------------------
% Theorem — Preservation of Intervals
% ---------------------------------------------------------

\begin{theorem}[Preservation of Intervals]
\label{thm:preservation-of-intervals}
Let $I$ be an interval and let $f : I \to \mathbb{R}$ be continuous
on $I$. Then $f(I)$ is an interval.

\smallskip
\noindent
\hyperref[prf:preservation-of-intervals]{\textit{Go to proof.}}
\end{theorem}

\begin{remark*}[Standard quantified statement]
\[
  I \text{ an interval},\;
  f \text{ cts on } I
  \;\Rightarrow\;
  f(I) \text{ is an interval.}
\]
\end{remark*}

\begin{remark*}[Interpretation]
This is the general form of the IVT, dropping the closed and bounded
hypotheses. The image of any interval under a continuous function is
an interval. This is the order-theoretic statement that continuity
preserves connectedness in $\mathbb{R}$: an interval has no gaps,
and a continuous image cannot introduce gaps that were not present in
the domain. The closed bounded case
(\hyperref[thm:image-of-closed-bounded-interval]{Theorem}) is a
special instance in which the image interval can be identified as
$[\min f, \max f]$.
\end{remark*}

% ---------------------------------------------------------
% Theorem — Darboux Property
% ---------------------------------------------------------

\begin{theorem}[Darboux Property]
\label{thm:darboux-property}
Let $f$ be continuous on $[a,b]$ and let $c \in \mathbb{R}$. If
$f(x) \neq c$ for all $x \in [a,b]$, then either $f(x) > c$ for
all $x \in [a,b]$, or $f(x) < c$ for all $x \in [a,b]$.

\smallskip
\noindent
\hyperref[prf:darboux-property]{\textit{Go to proof.}}
\end{theorem}

\begin{remark*}[Standard quantified statement]
\[
  f \text{ cts on } [a,b],\;
  \forall x \in [a,b]:\; f(x) \neq c
  \;\Rightarrow\;
  \bigl(\forall x \in [a,b]:\; f(x) > c\bigr)
  \;\vee\;
  \bigl(\forall x \in [a,b]:\; f(x) < c\bigr).
\]
\end{remark*}

\begin{remark*}[Interpretation]
The Darboux property is the contrapositive of the IVT: if the value
$c$ is not attained, then $f$ lies entirely on one side of $c$. A
continuous function on a closed interval cannot straddle a value
without hitting it. This is used in sign-argument proofs where one
needs to conclude that a nonzero function maintains its sign
throughout an interval.
\end{remark*}
